\section{南宋的重建：把失地后的国家安放在江海之间}

南渡后的朝廷并非把北宋原样搬到杭州。行在、临安、两浙与长江防线构成一个必须在战争中重新安放的国家：官员、宗室和军队需要俸粮，流民需要落脚，江南既有的田土、手工业、市场和港口则要在新的税役与军需中维持生产。一一二七年的政权断裂使许多家庭失去原有的籍贯、产业与关系；把南宋称为“偏安”，会遮蔽它在南方社会中重新组织财政、法律与日常生活的努力。

\subsection{和议不是停止治理}

绍兴和议为南宋争得重整时间，也把淮河与长江之间的防务、岁币和使节制度化。此后隆兴和议、开禧北伐及其失败，说明朝廷内部从未停止讨论恢复、守势和财政的相互关系。问题不在于哪一种口号更高尚，而在于远征所需的人、船、粮与信用能否持续供给；同样，守势也不是不作为，而是边堡、江防、军屯、转运和外交的长期投入。

这种体制有可见的收益。东南的税粮、漕运与商业能够支持一个失去北方核心地区的朝廷，科举、学校、司法与文书仍让中央和州县保持联系。它的代价也由此增长：军费和币制压力更集中地落在富庶地区，地方家庭面对的赋役、借贷和物价风险会随战事与灾害变化。不能把一纸和约、一个财政数字或一座港口的繁盛当作全国状态；南宋的力量恰是高度区域化的力量。

\subsection{海域、市场与迁徙}

海上贸易在南宋财政和消费中很重要，却不是与国家无关的自由市场。市舶、税收、航路安全、港口管理和商人网络相互纠缠；海外来货、沿海渔盐和内河转运又分别受不同季节、风向与地方秩序制约。纸币与会子扩大了支付和调度的可能，也使信用、折换和发行纪律成为政治问题。所谓“商业繁荣”应当包括这些制度摩擦，而不仅是把城市描绘为商品的乐园。

迁徙使南方社会的变化更为具体。北来官户和军户并不构成同质群体，本地居民、佃户、手工业者、僧侣和海商也不会只被动接受新的秩序。婚姻、宗族、寺院、市场和地方文书都是人们重建关系的场所。正因这些关系常在州县档案之外形成，国家材料所列的户口、田赋和军额既不可少，也不能当作社会的全部。

\subsection{从襄阳到海上：防线崩解的社会时间}

蒙古对四川、荆襄和江淮的推进，改变了南宋原有的防御节奏。襄阳、樊城长期被围，反映的是水陆通道、城防、运输和各方军事技术的竞争，而不只是某次城战的英雄叙事。城陷之后，地方仍有守城者、降附者、逃亡者和继续经营生计的人；一二七九年的崖山标志南宋政权终结，却不是江南社会一夜之间换了人间。

《宋史》本纪和《建炎以来系年要录》保留了朝廷决策与年序，《宋会要辑稿》有助于追问制度，笔记、墓志、契约、港口遗址和钱币则从不同角度接近区域生活。南宋末年的忠义叙事与元初的接收叙事都有强烈的立场。把它们并读，才能既不抹去战争中的选择，也不把后来形成的道德分类当成每个当时人的唯一动机。
